% Options for packages loaded elsewhere
\PassOptionsToPackage{unicode}{hyperref}
\PassOptionsToPackage{hyphens}{url}
\documentclass[
]{article}
\usepackage{xcolor}
\usepackage[margin=1in]{geometry}
\usepackage{amsmath,amssymb}
\setcounter{secnumdepth}{-\maxdimen} % remove section numbering
\usepackage{iftex}
\ifPDFTeX
  \usepackage[T1]{fontenc}
  \usepackage[utf8]{inputenc}
  \usepackage{textcomp} % provide euro and other symbols
\else % if luatex or xetex
  \usepackage{unicode-math} % this also loads fontspec
  \defaultfontfeatures{Scale=MatchLowercase}
  \defaultfontfeatures[\rmfamily]{Ligatures=TeX,Scale=1}
\fi
\usepackage{lmodern}
\ifPDFTeX\else
  % xetex/luatex font selection
\fi
% Use upquote if available, for straight quotes in verbatim environments
\IfFileExists{upquote.sty}{\usepackage{upquote}}{}
\IfFileExists{microtype.sty}{% use microtype if available
  \usepackage[]{microtype}
  \UseMicrotypeSet[protrusion]{basicmath} % disable protrusion for tt fonts
}{}
\makeatletter
\@ifundefined{KOMAClassName}{% if non-KOMA class
  \IfFileExists{parskip.sty}{%
    \usepackage{parskip}
  }{% else
    \setlength{\parindent}{0pt}
    \setlength{\parskip}{6pt plus 2pt minus 1pt}}
}{% if KOMA class
  \KOMAoptions{parskip=half}}
\makeatother
\usepackage{color}
\usepackage{fancyvrb}
\newcommand{\VerbBar}{|}
\newcommand{\VERB}{\Verb[commandchars=\\\{\}]}
\DefineVerbatimEnvironment{Highlighting}{Verbatim}{commandchars=\\\{\}}
% Add ',fontsize=\small' for more characters per line
\usepackage{framed}
\definecolor{shadecolor}{RGB}{248,248,248}
\newenvironment{Shaded}{\begin{snugshade}}{\end{snugshade}}
\newcommand{\AlertTok}[1]{\textcolor[rgb]{0.94,0.16,0.16}{#1}}
\newcommand{\AnnotationTok}[1]{\textcolor[rgb]{0.56,0.35,0.01}{\textbf{\textit{#1}}}}
\newcommand{\AttributeTok}[1]{\textcolor[rgb]{0.13,0.29,0.53}{#1}}
\newcommand{\BaseNTok}[1]{\textcolor[rgb]{0.00,0.00,0.81}{#1}}
\newcommand{\BuiltInTok}[1]{#1}
\newcommand{\CharTok}[1]{\textcolor[rgb]{0.31,0.60,0.02}{#1}}
\newcommand{\CommentTok}[1]{\textcolor[rgb]{0.56,0.35,0.01}{\textit{#1}}}
\newcommand{\CommentVarTok}[1]{\textcolor[rgb]{0.56,0.35,0.01}{\textbf{\textit{#1}}}}
\newcommand{\ConstantTok}[1]{\textcolor[rgb]{0.56,0.35,0.01}{#1}}
\newcommand{\ControlFlowTok}[1]{\textcolor[rgb]{0.13,0.29,0.53}{\textbf{#1}}}
\newcommand{\DataTypeTok}[1]{\textcolor[rgb]{0.13,0.29,0.53}{#1}}
\newcommand{\DecValTok}[1]{\textcolor[rgb]{0.00,0.00,0.81}{#1}}
\newcommand{\DocumentationTok}[1]{\textcolor[rgb]{0.56,0.35,0.01}{\textbf{\textit{#1}}}}
\newcommand{\ErrorTok}[1]{\textcolor[rgb]{0.64,0.00,0.00}{\textbf{#1}}}
\newcommand{\ExtensionTok}[1]{#1}
\newcommand{\FloatTok}[1]{\textcolor[rgb]{0.00,0.00,0.81}{#1}}
\newcommand{\FunctionTok}[1]{\textcolor[rgb]{0.13,0.29,0.53}{\textbf{#1}}}
\newcommand{\ImportTok}[1]{#1}
\newcommand{\InformationTok}[1]{\textcolor[rgb]{0.56,0.35,0.01}{\textbf{\textit{#1}}}}
\newcommand{\KeywordTok}[1]{\textcolor[rgb]{0.13,0.29,0.53}{\textbf{#1}}}
\newcommand{\NormalTok}[1]{#1}
\newcommand{\OperatorTok}[1]{\textcolor[rgb]{0.81,0.36,0.00}{\textbf{#1}}}
\newcommand{\OtherTok}[1]{\textcolor[rgb]{0.56,0.35,0.01}{#1}}
\newcommand{\PreprocessorTok}[1]{\textcolor[rgb]{0.56,0.35,0.01}{\textit{#1}}}
\newcommand{\RegionMarkerTok}[1]{#1}
\newcommand{\SpecialCharTok}[1]{\textcolor[rgb]{0.81,0.36,0.00}{\textbf{#1}}}
\newcommand{\SpecialStringTok}[1]{\textcolor[rgb]{0.31,0.60,0.02}{#1}}
\newcommand{\StringTok}[1]{\textcolor[rgb]{0.31,0.60,0.02}{#1}}
\newcommand{\VariableTok}[1]{\textcolor[rgb]{0.00,0.00,0.00}{#1}}
\newcommand{\VerbatimStringTok}[1]{\textcolor[rgb]{0.31,0.60,0.02}{#1}}
\newcommand{\WarningTok}[1]{\textcolor[rgb]{0.56,0.35,0.01}{\textbf{\textit{#1}}}}
\usepackage{graphicx}
\makeatletter
\newsavebox\pandoc@box
\newcommand*\pandocbounded[1]{% scales image to fit in text height/width
  \sbox\pandoc@box{#1}%
  \Gscale@div\@tempa{\textheight}{\dimexpr\ht\pandoc@box+\dp\pandoc@box\relax}%
  \Gscale@div\@tempb{\linewidth}{\wd\pandoc@box}%
  \ifdim\@tempb\p@<\@tempa\p@\let\@tempa\@tempb\fi% select the smaller of both
  \ifdim\@tempa\p@<\p@\scalebox{\@tempa}{\usebox\pandoc@box}%
  \else\usebox{\pandoc@box}%
  \fi%
}
% Set default figure placement to htbp
\def\fps@figure{htbp}
\makeatother
\setlength{\emergencystretch}{3em} % prevent overfull lines
\providecommand{\tightlist}{%
  \setlength{\itemsep}{0pt}\setlength{\parskip}{0pt}}
\usepackage{bookmark}
\IfFileExists{xurl.sty}{\usepackage{xurl}}{} % add URL line breaks if available
\urlstyle{same}
\hypersetup{
  pdftitle={Feature Embedding Methods},
  pdfauthor={Mr.Pardeshi},
  hidelinks,
  pdfcreator={LaTeX via pandoc}}

\title{Feature Embedding Methods}
\author{Mr.Pardeshi}
\date{2026-03-06}

\begin{document}
\maketitle

\subsection{R Markdown}\label{r-markdown}

This is an R Markdown document. Markdown is a simple formatting syntax
for authoring HTML, PDF, and MS Word documents. For more details on
using R Markdown see \url{http://rmarkdown.rstudio.com}.

When you click the \textbf{Knit} button a document will be generated
that includes both content as well as the output of any embedded R code
chunks within the document. You can embed an R code chunk like this:

\begin{Shaded}
\begin{Highlighting}[]
\CommentTok{\# ==============================================================================}
\CommentTok{\# Title: Contrast Encoding: Effect vs. Helmert}
\CommentTok{\# Dataset: PlantGrowth (Built{-}in R dataset)}
\CommentTok{\# Purpose: Demonstrate model fitting and interpretation using different encodings}
\CommentTok{\# ==============================================================================}

\CommentTok{\# 1. Load and prepare the dataset}
\FunctionTok{data}\NormalTok{(}\StringTok{"PlantGrowth"}\NormalTok{)}
\FunctionTok{View}\NormalTok{(PlantGrowth)}

\CommentTok{\# Let\textquotesingle{}s quickly look at the group means to help us interpret the models later}
\FunctionTok{cat}\NormalTok{(}\StringTok{"{-}{-}{-} Group Means {-}{-}{-}}\SpecialCharTok{\textbackslash{}n}\StringTok{"}\NormalTok{)}
\end{Highlighting}
\end{Shaded}

\begin{verbatim}
## --- Group Means ---
\end{verbatim}

\begin{Shaded}
\begin{Highlighting}[]
\NormalTok{group\_means }\OtherTok{\textless{}{-}} \FunctionTok{tapply}\NormalTok{(PlantGrowth}\SpecialCharTok{$}\NormalTok{weight, PlantGrowth}\SpecialCharTok{$}\NormalTok{group, mean)}
\FunctionTok{print}\NormalTok{(group\_means)}
\end{Highlighting}
\end{Shaded}

\begin{verbatim}
##  ctrl  trt1  trt2 
## 5.032 4.661 5.526
\end{verbatim}

\begin{Shaded}
\begin{Highlighting}[]
\FunctionTok{cat}\NormalTok{(}\StringTok{"Grand Mean:"}\NormalTok{, }\FunctionTok{mean}\NormalTok{(PlantGrowth}\SpecialCharTok{$}\NormalTok{weight), }\StringTok{"}\SpecialCharTok{\textbackslash{}n\textbackslash{}n}\StringTok{"}\NormalTok{)}
\end{Highlighting}
\end{Shaded}

\begin{verbatim}
## Grand Mean: 5.073
\end{verbatim}

\begin{Shaded}
\begin{Highlighting}[]
\CommentTok{\# ==============================================================================}
\CommentTok{\# 2. EFFECT ENCODING (Sum/Deviation Coding)}
\CommentTok{\# ==============================================================================}
\CommentTok{\# Concept: Compares the mean of a specific level to the GRAND MEAN of all levels.}
\CommentTok{\# Useful when you want to know "Does this specific group perform above or }
\CommentTok{\# below the overall average?" rather than comparing it to a control group.}

\CommentTok{\# Apply Effect Encoding}
\CommentTok{\# contr.sum(3) creates a matrix for 3 levels:}
\CommentTok{\# ctrl:  1  0}
\CommentTok{\# trt1:  0  1}
\CommentTok{\# trt2: {-}1 {-}1}
\FunctionTok{contrasts}\NormalTok{(PlantGrowth}\SpecialCharTok{$}\NormalTok{group) }\OtherTok{\textless{}{-}} \FunctionTok{contr.sum}\NormalTok{(}\DecValTok{3}\NormalTok{)}

\CommentTok{\# Fit the linear model}
\NormalTok{model\_effect }\OtherTok{\textless{}{-}} \FunctionTok{lm}\NormalTok{(weight }\SpecialCharTok{\textasciitilde{}}\NormalTok{ group, }\AttributeTok{data =}\NormalTok{ PlantGrowth)}

\FunctionTok{cat}\NormalTok{(}\StringTok{"========================================================}\SpecialCharTok{\textbackslash{}n}\StringTok{"}\NormalTok{)}
\end{Highlighting}
\end{Shaded}

\begin{verbatim}
## ========================================================
\end{verbatim}

\begin{Shaded}
\begin{Highlighting}[]
\FunctionTok{cat}\NormalTok{(}\StringTok{"SUMMARY: EFFECT ENCODING (SUM CONTRASTS)}\SpecialCharTok{\textbackslash{}n}\StringTok{"}\NormalTok{)}
\end{Highlighting}
\end{Shaded}

\begin{verbatim}
## SUMMARY: EFFECT ENCODING (SUM CONTRASTS)
\end{verbatim}

\begin{Shaded}
\begin{Highlighting}[]
\FunctionTok{cat}\NormalTok{(}\StringTok{"========================================================}\SpecialCharTok{\textbackslash{}n}\StringTok{"}\NormalTok{)}
\end{Highlighting}
\end{Shaded}

\begin{verbatim}
## ========================================================
\end{verbatim}

\begin{Shaded}
\begin{Highlighting}[]
\FunctionTok{print}\NormalTok{(}\FunctionTok{summary}\NormalTok{(model\_effect))}
\end{Highlighting}
\end{Shaded}

\begin{verbatim}
## 
## Call:
## lm(formula = weight ~ group, data = PlantGrowth)
## 
## Residuals:
##     Min      1Q  Median      3Q     Max 
## -1.0710 -0.4180 -0.0060  0.2627  1.3690 
## 
## Coefficients:
##             Estimate Std. Error t value Pr(>|t|)    
## (Intercept)   5.0730     0.1138  44.573   <2e-16 ***
## group1       -0.0410     0.1610  -0.255   0.8009    
## group2       -0.4120     0.1610  -2.560   0.0164 *  
## ---
## Signif. codes:  0 '***' 0.001 '**' 0.01 '*' 0.05 '.' 0.1 ' ' 1
## 
## Residual standard error: 0.6234 on 27 degrees of freedom
## Multiple R-squared:  0.2641, Adjusted R-squared:  0.2096 
## F-statistic: 4.846 on 2 and 27 DF,  p-value: 0.01591
\end{verbatim}

\begin{Shaded}
\begin{Highlighting}[]
\CommentTok{\# INTERPRETATION NOTES (Effect Encoding):}
\CommentTok{\# {-} Intercept: The Unweighted Grand Mean (average of the 3 group means).}
\CommentTok{\# {-} group1 (ctrl): The difference between the \textquotesingle{}ctrl\textquotesingle{} mean and the Grand Mean.}
\CommentTok{\# {-} group2 (trt1): The difference between the \textquotesingle{}trt1\textquotesingle{} mean and the Grand Mean.}
\CommentTok{\# {-} Note: trt2\textquotesingle{}s effect is not directly in the summary, but it can be calculated }
\CommentTok{\#   as {-}(group1 + group2) because the sum of deviations from the mean is zero.}


\CommentTok{\# ==============================================================================}
\CommentTok{\# 3. HELMERT ENCODING}
\CommentTok{\# ==============================================================================}
\CommentTok{\# Concept: Compares each level of a categorical variable to the mean of the }
\CommentTok{\# SUBSEQUENT (or previous) levels. }
\CommentTok{\# Useful in ordinal data or when testing sequential treatments (e.g., comparing }
\CommentTok{\# a new drug to the average of all previously existing drugs).}

\CommentTok{\# Apply Helmert Encoding}
\CommentTok{\# contr.helmert(3) creates a matrix:}
\CommentTok{\# ctrl: {-}1 {-}1}
\CommentTok{\# trt1:  1 {-}1}
\CommentTok{\# trt2:  0  2}
\FunctionTok{contrasts}\NormalTok{(PlantGrowth}\SpecialCharTok{$}\NormalTok{group) }\OtherTok{\textless{}{-}} \FunctionTok{contr.helmert}\NormalTok{(}\DecValTok{3}\NormalTok{)}

\CommentTok{\# Fit the linear model}
\NormalTok{model\_helmert }\OtherTok{\textless{}{-}} \FunctionTok{lm}\NormalTok{(weight }\SpecialCharTok{\textasciitilde{}}\NormalTok{ group, }\AttributeTok{data =}\NormalTok{ PlantGrowth)}

\FunctionTok{cat}\NormalTok{(}\StringTok{"}\SpecialCharTok{\textbackslash{}n}\StringTok{========================================================}\SpecialCharTok{\textbackslash{}n}\StringTok{"}\NormalTok{)}
\end{Highlighting}
\end{Shaded}

\begin{verbatim}
## 
## ========================================================
\end{verbatim}

\begin{Shaded}
\begin{Highlighting}[]
\FunctionTok{cat}\NormalTok{(}\StringTok{"SUMMARY: HELMERT ENCODING}\SpecialCharTok{\textbackslash{}n}\StringTok{"}\NormalTok{)}
\end{Highlighting}
\end{Shaded}

\begin{verbatim}
## SUMMARY: HELMERT ENCODING
\end{verbatim}

\begin{Shaded}
\begin{Highlighting}[]
\FunctionTok{cat}\NormalTok{(}\StringTok{"========================================================}\SpecialCharTok{\textbackslash{}n}\StringTok{"}\NormalTok{)}
\end{Highlighting}
\end{Shaded}

\begin{verbatim}
## ========================================================
\end{verbatim}

\begin{Shaded}
\begin{Highlighting}[]
\FunctionTok{print}\NormalTok{(}\FunctionTok{summary}\NormalTok{(model\_helmert))}
\end{Highlighting}
\end{Shaded}

\begin{verbatim}
## 
## Call:
## lm(formula = weight ~ group, data = PlantGrowth)
## 
## Residuals:
##     Min      1Q  Median      3Q     Max 
## -1.0710 -0.4180 -0.0060  0.2627  1.3690 
## 
## Coefficients:
##             Estimate Std. Error t value Pr(>|t|)    
## (Intercept)  5.07300    0.11381  44.573  < 2e-16 ***
## group1      -0.18550    0.13939  -1.331  0.19439    
## group2       0.22650    0.08048   2.814  0.00901 ** 
## ---
## Signif. codes:  0 '***' 0.001 '**' 0.01 '*' 0.05 '.' 0.1 ' ' 1
## 
## Residual standard error: 0.6234 on 27 degrees of freedom
## Multiple R-squared:  0.2641, Adjusted R-squared:  0.2096 
## F-statistic: 4.846 on 2 and 27 DF,  p-value: 0.01591
\end{verbatim}

\begin{Shaded}
\begin{Highlighting}[]
\CommentTok{\# INTERPRETATION NOTES (Helmert Encoding):}
\CommentTok{\# {-} Intercept: The Unweighted Grand Mean.}
\CommentTok{\# {-} group1: Compares Level 2 (trt1) to Level 1 (ctrl). }
\CommentTok{\#   Because the contrast values are 1 and {-}1 (difference of 2), the coefficient }
\CommentTok{\#   represents exactly HALF the difference between trt1 and ctrl means.}
\CommentTok{\# {-} group2: Compares Level 3 (trt2) to the AVERAGE of Level 1 \& 2 (ctrl \& trt1).}
\CommentTok{\#   Because the contrast values are 2 and {-}1 (difference of 3), the coefficient }
\CommentTok{\#   represents exactly ONE THIRD of the difference between trt2 and the average }
\CommentTok{\#   of the first two groups.}
\CommentTok{\#}
\CommentTok{\# *Note on modern Data Science tools: Some libraries normalize Helmert matrices }
\CommentTok{\# so the coefficients represent the exact differences rather than fractions. }
\CommentTok{\# R uses the unnormalized mathematical standard by default.}
\end{Highlighting}
\end{Shaded}


\end{document}
